\section{8.1 팀 프로젝트: 구조와 발표}
\begin{frame}{왜 지금 프로젝트를 하는가}
  \begin{itemize}
    \item Ch\,2$\sim$7: 밴딧 $\to$ n-step TD·Dyna-Q까지 표 기반 RL 이론을
      손에 쥨 --- Sutton \& Barto 표준 교과서와 거의 1:1.
    \item 지금까지 해본 것: \textbf{``교재가 고른 환경, 교재가 고른
      파라미터, 교재가 고른 시드''}에서 코드가 동작하는지 확인.
    \item 지금 필요한 것: 한 단계 위로 --- \textbf{환경을 고르고, 알고리즘
      두 개를 고르고, 결과가 흔들리면 시드와 프로토콜을 점검해서 다시
      고르는 것}까지.
    \item ``코드가 동작한다''는 증거 $\neq$ ``이 문제를 내가 설계하고
      평가할 수 있다''는 증거.
  \end{itemize}
  \vskip0.6em
  \begin{block}{프로젝트를 마치면 손에 넣는 것}
    보고서 + 발표 + peer-review --- ``설계하고 평가할 수 있다''의 증거.
  \end{block}
\end{frame}
